% Implementation chapter (30% review)
% Needs in the preamble: \usepackage{booktabs,listings,amsmath}

\chapter{Implementation}

This chapter describes the work completed for the 30\% review: the
development environment, profiling of the target workload on the Shakti
C-Class core, the design and RTL implementation of three custom
instructions (\texttt{FDOT2.S}, \texttt{PFMA.S} and \texttt{PRELU6.S}), the
verification method, and the measured results.

\section{Development Environment}

The baseline is the open-source Shakti C-Class core (IIT Madras), official
release 5.6.1. It is a 6-stage, in-order, 64-bit core implementing
RV64IMAFDC, written in Bluespec SystemVerilog (BSV). The tool flow is:

\begin{itemize}
  \item \textbf{Bluespec compiler (bsc)}: compiles the BSV source into Verilog.
  \item \textbf{Verilator}: builds a cycle-accurate simulator (\texttt{bin/out})
        from the generated Verilog.
  \item \textbf{RISC-V GCC 14} (\texttt{riscv64-unknown-elf-gcc}): compiles
        bare-metal C benchmarks; custom instructions are emitted with the
        assembler directive \texttt{.insn}.
  \item \textbf{elf2hex}: converts the program into a memory image loaded by
        the simulator.
\end{itemize}

Each simulation produces \texttt{rtl.dump}, a trace with one line per
retired instruction. Programs report pass or fail by writing to the
\texttt{tohost} address (\texttt{0x80001000}): \texttt{1} means pass. Timing
is measured with the hardware counters \texttt{mcycle} (clock cycles) and
\texttt{minstret} (instructions retired).

\section{Workload Profiling}

\subsection{Profiling on a PC}

The target workload is MobileNetV2 (feature extractor) followed by a GRU
(temporal model) over 8 video frames. A PyTorch profile per frame gave:

\begin{table}[h]
\centering
\begin{tabular}{lr}\toprule
Operation & MACs per frame\\\midrule
Pointwise $1\times1$ convolution & 267.9 M\\
Depthwise $3\times3$ convolution & 20.7 M\\
First $3\times3$ convolution & 10.8 M\\
GRU (8 steps) & 9.4 M\\\bottomrule
\end{tabular}
\caption{Multiply-accumulate (MAC) operations per frame.}
\end{table}

\subsection{Profiling on the Shakti core}

A PC profile does not show where an embedded core spends its time. Every
operation of the model was therefore written as a bare-metal C kernel, run
on the \emph{baseline} Shakti core with the real layer shapes, and timed with
\texttt{mcycle}. The measured cost per MAC (or per element) was multiplied by
the number of times each operation runs in one full inference.

\begin{table}[h]
\centering
\begin{tabular}{lrr}\toprule
Operation & Cycles per MAC & Share of time\\\midrule
Pointwise convolution & 11.7--13.9 & \textbf{83.8\%}\\
Depthwise convolution & 16.7 & 8.7\%\\
First convolution & 17.7 & 4.8\%\\
ReLU6 & 14.3 / element & 2.2\%\\
GRU & 12.0 & 0.4\%\\\midrule
\textbf{Total (one inference)} & & \textbf{31.9 billion cycles}\\\bottomrule
\end{tabular}
\caption{Baseline Shakti profile of MobileNetV2 + GRU.}
\end{table}

Pointwise convolution dominates, and every convolution costs 11--18 cycles
per MAC. The model is FP32, and Shakti already has a fused multiply-add
(\texttt{fmadd.s}), so a scalar custom MAC instruction would bring no
benefit. The instructions were therefore designed to do \emph{more work per
instruction} by packing two FP32 values into one 64-bit register.

\section{Instruction Encoding}

All new instructions use the RISC-V \textbf{custom-2} opcode
(\texttt{0x5B}), which is reserved for custom extensions, with the R4
(three-source-register) layout. The format bits \texttt{[26:25]} select the
instruction.

\begin{table}[h]
\centering
\begin{tabular}{llll}\toprule
fmt & Instruction & Internal FPU code & Assembly\\\midrule
\texttt{00} & \texttt{FDOT2.S}  & \texttt{0110} & \texttt{.insn r4 0x5B, 0, 0, fd, fs1, fs2, fs3}\\
\texttt{10} & \texttt{PFMA.S}   & \texttt{0111} & \texttt{.insn r4 0x5B, 0, 2, fd, fs1, fs2, fs3}\\
\texttt{01} & \texttt{PRELU6.S} & \texttt{0101} & \texttt{.insn r4 0x5B, 0, 1, fd, fs1, fs1, fs1}\\\bottomrule
\end{tabular}
\caption{Encoding of the custom instructions.}
\end{table}

\section{Design 1: FDOT2.S (Packed Dot Product)}

\[
fs1 = [a_1 \mid a_0],\quad fs2 = [b_1 \mid b_0],\quad
fd = a_1 b_1 + (a_0 b_0 + fs3)
\]

FDOT2 performs two multiply-adds in one instruction using the existing FP32
FMA unit twice in sequence. Because the rounding is identical to two
\texttt{fmadd.s} instructions, the result is bit-identical to normal code.

\subsection{Decoder remap}

Adding a completely new opcode would require changes across several decoder
functions. Instead, at the start of \texttt{fn\_decode}, any custom-2 word is
rewritten into an internal FMADD form with fmt = \texttt{10} (half
precision, which this core does not implement, so the pattern is otherwise
unused). The existing FMADD path then handles register reads, the
scoreboard and writeback unchanged. A separate 4-bit FPU code
(\texttt{fn}) tells the FPU which custom operation to perform.

\begin{lstlisting}[language=Verilog,basicstyle=\ttfamily\small]
// src/decoder.bsv
if (inst_in[6:0] == 'b1011011)                       // custom-2
  inst = {inst_in[31:27], 2'b10, inst_in[24:7], 7'b1000011};
...
if (inst_in[6:0] == 'b1011011)
  fn = (inst_in[26:25] == 2'b10) ? 4'b0111            // PFMA
     : ((inst_in[26:25] == 2'b01) ? 4'b0101 : 4'b0110); // PRELU6 / FDOT2
\end{lstlisting}

\subsection{Problems solved during FDOT2}

\begin{table}[h]
\centering\small
\begin{tabular}{p{4.3cm}p{4.3cm}p{5cm}}\toprule
Symptom & Cause & Fix\\\midrule
BSC rule conflict error & One rule both read the FMA response and issued a new request & Split the second step into its own rule\\
Illegal-instruction trap & Stale build: old decoder still in the simulator & Clean rebuild, check build time\\
Result was NaN & FPU NaN-boxing check replaced packed operands & Read the raw 64-bit operands\\
Still NaN & Operation was selected from bits that hold register data in R4 & Dedicated FPU code \texttt{fn = 0110}\\\bottomrule
\end{tabular}
\caption{Issues found and fixed while implementing FDOT2.}
\end{table}

\subsection{Key finding}

On a pointwise layer (16$\rightarrow$96 channels, 98,304 MACs), FDOT2 reduced
instructions by $1.87\times$ but cycles by only $1.07\times$. The FPU accepts
one operation at a time and each FMA takes about 10 cycles, so the two
sequential FMAs still waited for each other. The kernel is limited by
\emph{FMA throughput}, not instruction count. To go faster, two FMAs must run
\emph{in parallel}.

\section{Design 2: PFMA.S (Two-Lane Packed FMA)}

\[
fs1 = [x_1 \mid x_0],\quad fs2 = w,\quad fs3 = [c_1 \mid c_0],\quad
fd = [\,x_1 w + c_1 \mid x_0 w + c_0\,]
\]

PFMA computes two outputs at once. The two inputs are neighbouring pixels of
one channel, which sit next to each other in memory, so one 64-bit load
fetches both. Both lanes share the same weight. A \textbf{second FP32 FMA
unit} (\texttt{spfma2}) was added to the FPU so that both lanes execute in
parallel; each lane is exactly one \texttt{fmadd.s}, so the results are
bit-identical.

\begin{lstlisting}[language=Verilog,basicstyle=\ttfamily\small]
// src/fpu/hardfloat/fpu_hardfloat.bsv
let spfma2 <- mkspfma_instance;          // second FMA unit (lane 1)
...
if (opcode == 4'b0111) begin              // PFMA.S
  spfma.request (2'b00, r1[31:0],  r2[31:0], r3[31:0],  f3);
  spfma2.request(2'b00, r1[63:32], r2[31:0], r3[63:32], f3);
  rg_pfma <= True;                        // pack both results on output
end
\end{lstlisting}

The output rule joins the two 32-bit results into one 64-bit value and
combines their exception flags.

\subsection{Reuse for depthwise and first convolution}

Depthwise and the first convolution have the same structure (neighbouring
output pixels share the same weights), so PFMA was reused with only C code
changes. Because a 64-bit load must be 8-byte aligned, the input is first
copied into a ``pair'' layout. The cost of this copy is included in all
measurements.

\section{Design 3: PRELU6.S (Packed ReLU6)}

\[
fd = [\,\mathrm{clamp}_{[0,6]}(x_1) \mid \mathrm{clamp}_{[0,6]}(x_0)\,]
\]

ReLU6 needs no floating-point arithmetic. For IEEE-754 floats, positive values
order the same way as their bit patterns read as integers, and negative
values have the sign bit set. ReLU6 is therefore implemented with integer
comparisons and completes in \textbf{one cycle}:

\begin{lstlisting}[language=Verilog,basicstyle=\ttfamily\small]
else if (opcode == 4'b0101) begin   // PRELU6.S
  Bit#(32) q0 = (l0[31]==1 && l0[30:0]!=0) ? 0
              : ((l0[31]==0 && l0 > 32'h40C00000) ? 32'h40C00000 : l0);
  // same for lane 1 (q1)
  tx_fbox_out.u.enq(XBoxOutput{valid: True, data: {q1, q0}, fflags: 0});
end
\end{lstlisting}

Here \texttt{0x40C00000} is the bit pattern of 6.0.

\section{Summary of Hardware Changes}

\begin{table}[h]
\centering
\begin{tabular}{lp{9.5cm}}\toprule
File & Change\\\midrule
\texttt{decoder.defines} & Encodings of the custom and internal instruction forms\\
\texttt{decoder.bsv} & custom-2 remap to internal FMADD; FPU code selection\\
\texttt{fpu\_hardfloat.bsv} & Second FMA unit; FDOT2 two-step logic; PFMA lane packing; PRELU6 logic\\
\texttt{stage3.bsv} & Minor change in the execute stage\\\bottomrule
\end{tabular}
\caption{Files modified relative to Shakti C-Class 5.6.1.}
\end{table}

The rest of the pipeline (fetch, scoreboard, forwarding, memory and
writeback) is unchanged: the new instructions travel through the existing
floating-point path.

\section{Verification}

\begin{itemize}
  \item \textbf{Directed test}: FDOT2 on $[2, 1.5]\cdot[4, 3] + 0.5$ returned
        exactly $13.0$ (\texttt{0x41500000}).
  \item \textbf{Bit-exact benchmarks}: every benchmark runs the same
        computation with standard RISC-V code and with the custom instruction,
        then compares all outputs bit for bit and reports through
        \texttt{tohost}. All benchmarks pass.
  \item \textbf{Regression}: the official \texttt{rv64uf/fmadd.S} test from
        riscv-tests passes after every hardware change, confirming that the
        standard FMADD instruction is unaffected.
\end{itemize}

\section{Results}

All results below were measured on the simulated Shakti core with
\texttt{mcycle}, using real MobileNetV2 layer shapes, and are bit-exact.

\begin{table}[h]
\centering
\begin{tabular}{lccc}\toprule
Kernel & Baseline & Custom & Speedup\\\midrule
Pointwise, $K$ = 16--144 & 11.1--11.8 cyc/MAC & 6.6--7.0 & \textbf{1.69--1.70$\times$}\\
Pointwise, $K$ = 320--960 & 13.8--13.9 cyc/MAC & 9.2--9.3 & \textbf{1.49--1.50$\times$}\\
Depthwise $3\times3$ (incl.\ copy) & 16.99 cyc/MAC & 10.86 & \textbf{1.56$\times$}\\
First conv $3\times3$, stride 2 & 18.28 cyc/MAC & 9.22 & \textbf{1.98$\times$}\\
ReLU6 (PRELU6) & 14.00 cyc/elem & 3.52 & \textbf{3.98$\times$}\\\bottomrule
\end{tabular}
\caption{Measured per-kernel results ($K$ = number of input channels).}
\end{table}

Large pointwise layers gain less because their inputs and weights
(30--90\,KB) do not fit in the 16\,KB data cache; PFMA speeds up the
arithmetic but not the memory stalls.

\begin{table}[h]
\centering
\begin{tabular}{lrr}\toprule
Configuration & Cycles & Speedup\\\midrule
Baseline Shakti & 31.9 B & 1.00$\times$\\
+ PFMA (pointwise, depthwise, first conv) & 19.8 B & 1.61$\times$\\
\textbf{+ PRELU6} & \textbf{19.3 B} & \textbf{1.66$\times$}\\\bottomrule
\end{tabular}
\caption{Whole-model estimate for one inference (8 frames + GRU).}
\end{table}

The whole-model figure is an estimate: the measured cycles per operation are
multiplied by the exact operation counts of the model, since simulating the
full network cycle by cycle would take days.

\section{Current Limitation}

The overall gain is limited by Amdahl's law. Pointwise convolution is still
about 85\% of the time, and the FPU is not pipelined: it accepts only one
operation at a time, so each PFMA waits about 10 cycles. Pipelining the FPU,
so that a new FMA can start every cycle, is the next major step, followed by
INT8 instructions for the TinyMetro-Net model.
